\documentclass[11pt]{article}
\usepackage[margin=1in]{geometry}
\usepackage{amsmath,amssymb,amsthm}
\usepackage[hidelinks]{hyperref}
\newtheorem{theorem}{Theorem}
\begin{document}
\setlength{\emergencystretch}{1em}
\title{Two remaining families for ordinary degree-seven binomial SFC}
\author{}\date{}\maketitle
Let $\mathcal L(X^j)=j!$ and put
\[
 P_j^{a,d}(t)=\sum_{k=0}^j\binom jk(aj+1)_{dk}t^k.
\]
For a nonzero phase $f$, Strong Factorial means that
$\mathcal L(f^m)$ and $\mathcal L(f^{m+1})$ cannot both vanish
at any positive integer $m$.

\begin{theorem}[Exact two-family frontier]\label{thm:two-family-frontier}
Strong Factorial for every ordinary complex univariate binomial
of degree at most seven is equivalent to
\[
 \gcd_{\mathbb Q[t]}(P_{N-1}^{a,d},P_N^{a,d})=1
\]
at every terminal integer $N\ge6$ in exactly the following two rows:
\[
 \begin{array}{c|c|l}
 (a,d)&\text{support}&\text{terminal restriction}\\ \hline
 (1,6)&\{1,7\}&N\notin\{b7^r:1\le b\le6,\ r\ge0\}\\
 (5,2)&\{5,7\}&N\equiv12,14\pmod{16}.
 \end{array}
\]
Support $\{3,7\}$ is proved at every positive pair.
For the first row, every hypothetical common root $t$ necessarily
satisfies $v_7(-1/t)=1$ at every embedding into $\mathbb C_7$.
For the second row, its gcd is either $1$ or a monic quadratic
\[
 t^2-\alpha t+\delta,\qquad
 \alpha,\delta\in\mathbb Q_{>0},\quad \alpha^2<4\delta.
\]
Thus excluding these quadratic factors is equivalent to completing
the second row. Coprimality on the two remaining rows is not asserted.
\end{theorem}
\begin{proof}
Theorem \texttt{thm:residue-frontier} of the
\href{ORDINARY_DEGREE_SEVEN_RESIDUE_FRONTIER.tex}{three-row residue classification}
reduces degree seven to support $\{1,7\}$ at every $N\ge6$,
support $\{3,7\}$ on three terminal progressions, and support
$\{5,7\}$ on $N\equiv12,14\pmod{16}$.
Theorem \texttt{thm:coprime} of the
\href{LOWER_THREE_FOURTH_POWER_FROBENIUS_SFC.tex}{fourth-power Frobenius-minor proof}
completes the entire middle support, including its three former
progressions.

Corollary \texttt{cor:pair} of the
\href{PHYSICAL_PRIME_TOTAL_DEGREE_REDUCTION.tex}{physical prime-total-degree reduction},
with $p=7$, proves the first support at each displayed excluded index
$N=b7^r$. Its Theorem \texttt{thm:radius} confines any multiple
root of $T_{N,N,6}$ to $v_7(-1/t)=1$.
The same physical reduction's coefficientwise consecutive-moment
identities show that a common zero of $P_{N-1}^{1,6},P_N^{1,6}$
would be a multiple root of $T_{N,N,6}$.
These statements hold over every algebraic embedding, since all the
polynomials have rational coefficients.

For the last row, Corollary \texttt{cor:factor-sign} of the
\href{LOWER_FIVE_QUADRATIC_HALF_PLANE_REDUCTION.tex}{quadratic half-plane reduction}
gives exactly the stated alternatives for the gcd at lower index
$n=N-1$. Its conclusion includes the all-index degree bound of the
earlier residue note.

The two remaining rows are therefore sufficient for the earlier
three-row criterion. Conversely they are subsets of the degree-seven
assertion, so they are necessary. This is an exact equivalence
without any finite cutoff or assertion about untested residues.
As in the earlier classification,
\[
 \mathcal L((X^a(A+BX^d))^j)=A^j(aj)!P_j^{a,d}(B/A)
 \qquad(AB\ne0).
\]
A common complex zero is equivalent to a nonconstant rational gcd;
nonzero degenerate monomials have all positive moments nonzero.
The earlier classification already supplies degree at most six and
all supports through terminal order five.
\end{proof}
The written Frobenius-minor argument closes support $\{3,7\}$.
The later \href{LOWER_FIVE_QUADRATIC_SECTOR_SFC.tex}{the sector theorem} excludes the $\{5,7\}$ quadratic factors
for every positive Gamma shape. The current degree-seven problem is the
remaining lower-one gap-six case described by \href{PHYSICAL_PRIME_ANALYTIC_INTERPOLATION.tex}{the later gap-six theorem}.
All-degree binomial SFC and unrestricted FC/SFC remain open. All review and replay here are internal;
no independent review, formal verification or priority is asserted.
\end{document}
